\section{Exact determinant multiplicities along the stretched Hessian family}

\subsection*{Finding O-028: invariant rectangle shifts of two-row cores}
Fix \(n\geq2\), put \(R=n^2\), and let
\[
 \Phi\in\mathbb C[\overline{\operatorname{GL}_R\cdot\det_n}]
\]
be a nonzero homogeneous \(\operatorname{SL}_R\)-invariant of
positive degree \(e\).
Write its polynomial highest weight as
\(\delta\mathbf1_R\), where \(\delta\in\mathbb Z_{>0}\).
Scalar weights give
\[
                              e=n\delta.
                                                               \tag{O-028.1}
\]
Such an invariant exists by the known polystability and
fundamental-invariant theorem used in M-020.
No numerical value of its degree is required below.

Let \(\sigma\vdash na\) have at most two rows, allowing the
empty partition when \(a=0\). For every integer \(j\geq0\), put
\[
               k=j\delta,\qquad
               \tau=\sigma+k\mathbf1_R,\qquad
               d=a+nk.
                                                               \tag{O-028.2}
\]
Then all the following actual and orbit multiplicities agree:
\[
 \boxed{\quad
 K_n(\tau)
   =\operatorname{sk}\bigl(\tau,(d^n)\bigr)
   =g\bigl(\tau,(d^n),(d^n)\bigr)
   =g\bigl(\sigma,(a^n),(a^n)\bigr)
   =K_n(\sigma).
 \quad}                                                       \tag{O-028.3}
\]
The two-row value on the right has the positive partition
difference formula in O-006 and F-002.
The assertion is about an actual rectangular semi-invariant
present in the determinant coordinate ring. An arbitrary
rectangle shift is not substituted for that hypothesis.

\begin{proof}[Preserving every two-row multiplicity copy]
Let \(A=\mathbb C[\overline{\operatorname{GL}_R\cdot\det_n}]\).
The determinant orbit closure is irreducible, so \(A\) is a
domain. Multiplication by \(\Phi^j\) therefore injects the
entire highest-weight space of weight \(\sigma\) into that
of weight \(\tau\), with degree changing from \(a\) to
\[
                         a+je=a+nk=d.
\]
Consequently \(K_n(\tau)\geq K_n(\sigma)\).

For the reverse bound, write \(W=U\otimes V\), where
\(\dim U=\dim V=n\). The usual determinant twist identity gives
\[
 \mathbb S_{\sigma+k\mathbf1_R}(U\otimes V)
   =\det(U\otimes V)^k\otimes\mathbb S_\sigma(U\otimes V).
\]
Since
\[
 \det(U\otimes V)^k=(\det U)^{nk}(\det V)^{nk},
\]
the multiplicity of
\((\det U)^d\otimes(\det V)^d\) on the left is exactly the
multiplicity of
\((\det U)^a\otimes(\det V)^a\) in
\(\mathbb S_\sigma(U\otimes V)\). Thus
\[
                 g(\tau,(d^n),(d^n))
                      =g(\sigma,(a^n),(a^n)).
                                                               \tag{O-028.4}
\]
This is the exact rectangular Kronecker shift identity;
it is an ordinary coefficient identity and uses no
assumed equality of orbit and closure multiplicities.

O-006 proves \(K_n(\sigma)=g(\sigma,(a^n),(a^n))\) for
every two-row \(\sigma\). The established comparison
\(K_n\leq\operatorname{sk}\leq g\), together with the
injection and (O-028.4), now bounds every term in
(O-028.3) between the same number. This proves equality.
The empty partition and \(j=0\) cases are included:
the degree-zero highest-weight space consists of constants.
\end{proof}

\paragraph{Application to the Hessian-bracket obstruction.}
For \(n\geq3\), use the A-030 notation
\[
 R=n^2,\qquad D=R+1,\qquad
 \lambda_n=(n+R(n-2)+1,\ 3,\ 2^{R-2}).
                                                               \tag{O-028.5}
\]
Set
\[
       t=(n-1)^2=D-2n,\qquad E=nt,\qquad
       h_n(s)=[z^s]\prod_{q=2}^{n}(1-z^q)^{-1}.
                                                               \tag{O-028.6}
\]
For every integer \(s\geq1\) satisfying \(\delta\mid2s\),
\[
 \boxed{\quad
       K_n(s\lambda_n)
       =\operatorname{sk}\bigl(s\lambda_n,((sD)^n)\bigr)
       =g\bigl(s\lambda_n,((sD)^n),((sD)^n)\bigr)
       =h_n(s).
 \quad}                                                       \tag{O-028.7}
\]
In particular, put
\[
                            p=\frac{\delta}{\gcd(\delta,2)}.
\]
The actual determinant multiplicity is known exactly for
every positive \(s\) divisible by \(p\). It is strictly
positive on this progression whenever \(s\geq2\).
The integer \(p\) is determined by the chosen actual
invariant; a numerical value is not asserted for general \(n\).

\begin{proof}
Padding all partitions with zeros to length \(R\), direct
subtraction gives
\[
        \lambda_n-2\mathbf1_R=(E-1,1),\qquad
        s\lambda_n=(s(E-1),s)+2s\mathbf1_R.
                                                               \tag{O-028.8}
\]
The core \(\sigma_s=(s(E-1),s)\) has size \(nst\), so
its outer degree is \(st\). The exact two-row formula
O-002, with rectangle height \(n\) and width \(st\), yields
\[
                    K_n(\sigma_s)
                   =g(\sigma_s,((st)^n),((st)^n))
                   =h_n(s).
                                                               \tag{O-028.9}
\]
The width bound is inactive because \(st\geq s\).
If \(\delta\mid2s\), multiply the complete core space by
\(\Phi^{2s/\delta}\). Its degree increases by
\[
                     e\,\frac{2s}{\delta}=2sn,
\]
and the total degree is \(st+2sn=sD\), exactly the degree
required for \(s\lambda_n\). Apply (O-028.3) to obtain
(O-028.7).

Finally \(h_n(1)=0\). For \(n\geq3\), both parts two and
three occur in (O-028.6), so every integer \(s\geq2\)
has a representation by those parts and \(h_n(s)>0\).
This proves the positivity claim on the stated progression.
\end{proof}

\paragraph{Orbit multiplicity at every dilation.}
The ordinary and symmetric orbit coefficients in fact satisfy
\[
 \operatorname{sk}\bigl(s\lambda_n,((sD)^n)\bigr)
   =g\bigl(s\lambda_n,((sD)^n),((sD)^n)\bigr)
   =h_n(s)\qquad(s\geq1),
                                                               \tag{O-028.9a}
\]
without the divisibility condition \(\delta\mid2s\).
For the ordinary coefficient this is the determinant twist
in (O-028.8)--(O-028.9).
For the symmetric coefficient, transpose on
\(\mathbb C^n\otimes\mathbb C^n\) has determinant
\((-1)^{n(n-1)/2}\).
The determinant-line twist of exponent \(2s\) therefore
has transpose scalar one.
Consequently the same identification of rectangular
multiplicity spaces preserves the transpose involution
and its symmetric subspace.
Its two-row core has symmetric coefficient equal to the
ordinary coefficient by O-006, proving (O-028.9a).
Thus the orbit multiplicity is already positive at every
\(s\geq2\); an actual invariant section is still needed
for the closure equalities proved above and below.

\paragraph{Every sufficiently large dilation, and a genuine support hole.}
An additional known theorem on the invariant degree monoid
strengthens the progression conclusion to all sufficiently
large integers. For every fixed \(n\geq3\), there is an integer
\(s_0(n)\) such that
\[
                       K_n(s\lambda_n)=h_n(s)
                  \qquad\text{for every }s\geq s_0(n).
                                                               \tag{O-028.11}
\]
No numerical bound on \(s_0(n)\) is asserted.
The exact ordinary and symmetric coefficients in
(O-028.7) equal the same value for all these \(s\).
In particular the actual determinant support monoid
\[
                        S_n=\{\mu:K_n(\mu)>0\}
\]
has the explicit relative-saturation hole
\[
                           \lambda_n\in
                         \operatorname{Sat}(S_n)\setminus S_n.
                                                               \tag{O-028.12}
\]
Here relative saturation uses the group generated by \(S_n\);
membership in that group is proved below, rather than inferred
from the existence of a single positive stretch.

\begin{proof}
Let \(\mathcal E_n\) be the monoid of degrees of nonzero
\(\operatorname{SL}_{n^2}\)-invariants in the determinant
coordinate ring.
The additional primary inputs are
B\"urgisser--Ikenmeyer, \emph{Fundamental invariants of orbit
closures}, arXiv:1511.02927v2:
\begin{itemize}
\item Theorem 2.5, printed p.~5, gives the determinant
stabilizer period
\[
 \alpha_n=
 \begin{cases}
  1,&n\bmod4\in\{0,1\},\\
  2,&n\bmod4\in\{2,3\};
 \end{cases}
\]
\item Definition 3.1 and Theorem 3.4, printed p.~8, say
that the invariant degree monoid generates the group
\[
                          \operatorname{gp}(\mathcal E_n)
                                    =n\alpha_n\mathbb Z.
\]
\end{itemize}
The theorem's polystability hypothesis is supplied by
Corollary 2.9, pp.~6--7, as already used in M-020.
Thus
\[
                    M_n=\frac1{n\alpha_n}\mathcal E_n
\]
is an additive submonoid of \(\mathbb Z_{\geq0}\) whose
generated group is \(\mathbb Z\).

For completeness, this condition implies that \(M_n\)
contains all sufficiently large integers.
Choose any positive \(h\in M_n\).
The residues of \(M_n\) in the finite cyclic group
\(\mathbb Z/h\mathbb Z\) form an additive submonoid,
and hence a subgroup: every element has its inverse
among its nonnegative multiples.
Since \(M_n\) generates \(\mathbb Z\), this subgroup is
all of \(\mathbb Z/h\mathbb Z\).
Choose \(b_j\in M_n\) for each residue \(j\).
Every integer \(u\geq\max_j b_j\), in its corresponding
residue class, is \(b_j+vh\) for some \(v\geq0\) and
therefore belongs to \(M_n\).
Let \(c_n\) be any bound with this property.

Because \(\alpha_n\) is one or two, \(2s/\alpha_n\)
is an integer for every \(s\).
If
\[
                  s\geq
             \max\{2,\lceil\alpha_n c_n/2\rceil\},
\]
then \(2s/\alpha_n\in M_n\).
Hence the determinant coordinate ring contains a
nonzero invariant \(\Psi_s\) of degree \(2sn\), whose
highest weight is \(2s\mathbf1_{n^2}\).
Apply the invariant-shift lemma (O-028.3) to
\(\Psi_s\) and the same two-row core
\(\sigma_s=(s(E-1),s)\).
It gives (O-028.11), with strictly positive value because
\(s\geq2\).
There is no need to express \(\Psi_s\) as a power of one
fixed fundamental invariant.

Choose \(s\) large enough that both \(s\lambda_n\) and
\((s+1)\lambda_n\) have positive determinant multiplicity.
Their difference shows
\[
                   \lambda_n=(s+1)\lambda_n-s\lambda_n
                                  \in S_n-S_n.
\]
Some positive multiple of \(\lambda_n\) is in \(S_n\),
so this proved group membership implies
\(\lambda_n\in\operatorname{Sat}(S_n)\).
Finally A-030 gives \(K_n(\lambda_n)=0\).
This proves the hole assertion (O-028.12).
\end{proof}

A-030 also proves that the actual unpadded permanent
multiplicity at \(s\lambda_n\) is positive for every \(s\geq1\),
by powers of its nonzero Hessian-bracket highest-weight section.
Thus the occurrence obstruction present at \(s=1\) disappears
for every sufficiently large dilation, with the exact
determinant multiplicity given by (O-028.11).
The argument uses the special two-row core and the
determinant's invariant period; it is not a general
Kronecker or determinant-support saturation theorem.

The current status is stronger than this invariant-based
argument alone. A-034 proves, for every \(n\geq3\) and
every integer \(s\geq1\),
\[
 q^n_{s\lambda_n}(s(n^2+1)[n])
       \geq p_n(s)>h_n(s)=K_n(s\lambda_n),
 \qquad
 p_n(s)=[z^s]\prod_{i=1}^{n}(1-z^i)^{-1}.
\]
The determinant equality follows by multiplying the
entire binary highest-weight space by a nonzero
Hessian section; it requires no invariant-period
threshold. Hence both actual multiplicities are
positive at every \(s\geq2\), and the strict gap
persists at all dilations. O-030 extends that
construction to a three-parameter flat-tail family.
For \(n=3\), F-028's injective binary
permanent--Hessian coefficient map supplies a stronger
permanent lower bound. The invariant-shift lemma and
numerical-monoid proof above remain valid independent
tools; their ray consequences are superseded by these
all-integer formulas.

\paragraph{A concrete even-size specialization.}
The known invariant computation for \(n=4\) makes one
progression numerical.
B\"urgisser--Ikenmeyer, \emph{Fundamental invariants of orbit
closures}, arXiv:1511.02927v2, Theorem 3.15,
printed p.~13, gives the general minimal-degree bound.
The paragraph following Proposition 3.28, printed p.~18,
reports the computation
\[
                            P_{4,16}(\det_4)\ne0.
\]
Together these statements give a nonzero invariant of
degree \(e=16\), and hence \(\delta=4\).
This nonvanishing is a used published computation, not
recomputed or relabeled as new here.
The original PDF is
\url{https://arxiv.org/pdf/1511.02927v2}, retained locally
as \path{external_sources/BI_fundamental_invariants.pdf}.

Here \(D=17\) and
\(\lambda_4=(37,3,2^{14})\).
For every integer \(u\geq1\), (O-028.7) becomes
\[
 \begin{split}
 K_4\bigl(2u(37,3,2^{14})\bigr)
   &=[z^{2u}]\frac1{(1-z^2)(1-z^3)(1-z^4)}\\
   &=\sum_{v=0}^{\lfloor u/3\rfloor}
                    \left(\left\lfloor\frac{u-3v}{2}\right\rfloor+1\right).
 \end{split}                                                   \tag{O-028.10}
\]
For the second equality, a solution to
\(2a+3b+4c=2u\) has \(b=2v\);
then \(a+3v+2c=u\), leaving exactly the displayed
number of choices for \(c\).
The first values for \(u=1,2,3,4,5,6\) are
\[
                              1,\ 2,\ 3,\ 4,\ 5,\ 7.
\]
These are actual determinant-closure multiplicities,
not merely ordinary rectangular coefficients.

\paragraph{Grothendieck dictionary and the integral limitation.}
The proof uses multiplication by a genuine nonzero section
of a one-dimensional representation, followed by the
actual coordinate-ring upper bound.
The determinant-twist identity is the Schur-functor
identity obtained by tensoring with a determinant line;
the special \(\lambda\)-structure is SGA~6, VI,
Theorem~3.3, pp.~392--393, with the functorial decomposition
in Macdonald, I, Appendix A (5.3)--(5.6), pp.~155--156.
Highest-weight products and rectangular characters are
the coordinate-comodule and grading constructions of
SGA~3, I, Propositions~4.7.2--4.7.3, pp.~34--36.
These tensor and determinant identities have integral
models. The actual section \(\Phi\), its nonvanishing on
the determinant, the irreducible multiplicity comparison,
and the exact assertions above are made over
\(\mathbb C\).
No chosen integral normalization of \(\Phi\) is assumed
nonzero on every special fiber.
The result uses neither normality nor an inference from
positivity in a virtual Grothendieck ring.
